Rounded vehicle clearance, optical occlusion, and pitch bounds jointly determine T1's observable region. Body clearance at $B_2$ limits the minimum-range arc, and the readiness tie-break selects V4 with readiness time \BReadyTime{}~s. V3 and V5 attain T2's analytical minimum local RMS of \BFirstRMS{}~m at an orthogonal intersection under the 30~s travel ceiling.

For concurrent alerts, beam search at $k=8$ matches the exact discrete latest start of \BBeamEight{}~s. Branch-and-bound checks all 1,000 vehicle pairings and 4,200 acyclic service orders on the 585-station pool. In the same discrete problem, enforcing a $30^\circ$ minimum crossing angle increases latest start by \BAngleDelay{}~s and lowers the selected exact witness's worst local RMS from \BAngleUnrestrictedRms{} to \BAngleGuardedRms{}~m. This makes the latency cost of a minimum-angle requirement explicit.

Under the Problem 3 assumptions, the 95\% pointing-coverage limit is \BIntegratedCrossing{}~m. At 80~m, the condition $a+80b\le\sigma^*$ gives a direct specification for the two drift parameters. For nominal hardware, the quadrature and Monte Carlo checks put 80~m below the coverage requirement.

In the 115 distinct stress cases, rescheduling reaches every alert. The V5 failure shows that interrupting a slew can increase the final start time even when completion stays the same. The zero-threshold selector retains Keep Plan for that case, and thresholds of 0--5~s give the same selected plans on the tested grid. Caching the nominal assignments also cuts latest start by 22.69~s in the pure 25~s outage case. Among the configurations tested at minimum latency, a shortlist of five stations reduces travel while a shortlist of eight raises mean observation quality.
